%!TEX root = ../thesis.tex
%*******************************************************************************
%*********************************** First Chapter *****************************
%*******************************************************************************

\chapter{State of the Art} 
\label{ch1:sec:sota}

\ifpdf
    \graphicspath{{State-Of-The-Art/Figs/Raster/}{State-Of-The-Art/Figs/PDF/}{State-Of-The-Art/Figs/}}
\else
    \graphicspath{{State-Of-The-Art/Figs/Vector/}{State-Of-The-Art/Figs/}}
\fi

In this chapter, I will present the state of the art of cultural competence in social robotics. Next, I will devote a section to cultural competence in ML. Finally, I will conclude by discussing cultural competence across robot functional domains.

%\color{red}
%I reference (nome anno) sono eleganti, ma sono talmente tanti che rendono il testo molto faticoso. Riesci a farli diventare [numerici]?\color{black}

%\color{red}
\section{Cultural factors in robotics}
%\color{black}

The integration of cultural factors in robotic systems has been a subject of rigorous investigation over the past decade \citep{Bartneck2007, Rehm2013}. 
Existing research has predominantly focused on identifying determinants of system acceptability across diverse cultural groups, considering either morphological appearance or behavioral patterns \citep{Shibata2009, Dingjun2010, Lee2016, vernon2024african, papadopoulos2022transcultural, recchiuto2019designing, akinade2023culturally, ornelas2023redefining, zantou2023culturally, Papadopoulos_CARESSES, battistuzzi2021socially, papadopoulos2020caresses, 8673086}. 
This includes the study of cultural nuances in verbal and non-verbal communication \citep{Rau2009, Andrist2015, Bruno2018, Rehm2018, CloudCultureDialogue, bruno2016fuzzy, chugh2024aisles, obremski2022impact, nias2020culturebot, recchiuto2020feasibility, GreetingSelection, pena2019eeva}, as well as variations in perceived social distance and proxemics \citep{Eresha2013, Joosse2014, Trovato2015, truong2016towards}.

Beyond social cues, the literature has addressed cultural competence within specific AI sub-domains, such as automated planning \citep{Planning, fedriga2023linguistic}, knowledge representation \citep{GreetingSelection, KnowledgeRepresentation, Grassi_2022, mansouri2020can, bringsjord2020culturally}, and activity recognition \citep{Chiang2019}. 
However, as evidenced by recent systematic and scoping reviews in social robotics for geriatric assistance, these approaches have generally failed to establish a unified conceptual framework. 
There remains a critical lack of generalized architectures designed to embed comprehensive cultural competence into intelligent systems for health and social care \citep{Abdi2018, Pu2019, Maresova2020, Lee2020}.

Recently, a conceptual framework was proposed\footnote{https://cordis.europa.eu/article/id/124441-the-worlds-first-culturally-sensitive-robots-for-elderly-care} to develop socially assistive robots (SARs) for geriatric care with cultural competence \citep{PavingTheWayCulturallyCompetent, KnowledgeRepresentation, Papadopoulos_CARESSES}. 
This framework was developed by synthesizing foundational research from transcultural nursing \citep{Leininger1988} and established paradigms in culturally competent healthcare \citep{Betancourt2003}.

%\color{red}
\section{ML and cultural factors}
%\color{black}
Anthropological studies of material culture demonstrate that everyday objects and the domestic interiors in which they are situated exhibit systematic design variations across different global regions, even when such objects share equivalent functional properties. This variation is expected, as object design is closely linked to the social functions it supports and to the regional availability of materials used in fabrication \citep{Robb2015}. Consequently, while certain aesthetic universals have been identified \citep{Hardonk1999, Hekkert2008}, specific visual features tend to recur more frequently and coherently within particular geographical and cultural domains.

Metaphorically, two objects may share identical functional affordances \citep{Gibson1986} yet still exhibit distinct “cultural signatures”. These features constitute relevant information that machine learning algorithms should account for during classification in order to ensure robustness. However, existing computational approaches have largely overlooked these nuances, thereby neglecting the role of cultural factors in improving model performance.

From a computational perspective, recent research has investigated principles and architectures for the design of culture-aware machine learning models \citep{miehle2021multimodal, khan2021one, dmytryk2022mining, karpouzis2024plato}. Significant emphasis has been placed on the development and evaluation of natural language processing techniques aimed at improving cultural sensitivity \citep{blanchard2024cultural, bhatia2023gd, almaliki2016persuasive, dlHumorRecognition, albuquerque2024hype, Wu2012TheAO, cirucci2024culturally, kuo2024identification, dwivedi2024navigating}.

Additional research directions include the integration of cultural considerations into generative image tasks \citep{Quan2018, Zhou2022445, yun2024cic, mohamed2022artelingo} and classification problems across diverse application domains. These range from analyzing vaccine hesitancy \citep{alharbi2023cultural, alharbi2024deep}, to studying digital discourse on mental health during the COVID-19 pandemic \citep{castilla2022applying}, as well as AI-augmented clinical triage \citep{jordan2023impact} and personalized recommender systems \citep{yohe2023towards}.

%\color{red} Mi chiedo se mettere questo. Lo hai già detto bene nell'introduzione. Qui limitati a parlare dello stato dell'arte. 
%\color{blue}
%Rimosso.
%\color{black}


ML has achieved unprecedented performance benchmarks, frequently exceeding human-level proficiency across diverse domains \citep{openai2023gpt4, alphago, minerva, human_performance}. 
These advancements have catalyzed researchers and practitioners across various disciplines to leverage ML in addressing domain-specific challenges. 
This paradigm shift has resulted in a significant proliferation of successful ML applications, establishing the technology as a ubiquitous force. 
Consequently, ML continues to revolutionize and profoundly influence the broader societal landscape \citep{MLForSociety, MLTrends}.

ML plays a pivotal role in robotics, significantly enhancing the capacity of autonomous agents to perceive and interact with complex environments \citep{DLRobotics, RLRobotics, MPRobotics, xu2020learning, bredeche2006perceptual, nguyen2020deep, zhou2022computer, arzani2020switching, hong2020multimodal, gao2021hybrid, yu2021deep}. 
Within the social robotics domain, ML is equally critical, as it provides the computational frameworks necessary for robots to interpret and model human environments. 
These capabilities enable agents to exploit environmental and social data to optimize interactions with human users across both verbal and non-verbal modalities \citep{SocialRobotics, zheng2013designing, jacob2015optimal, saunderson2020investigating, ogenyi2019physical, rakovic2022gaze, alsamhi2022autonomous}.

The most advanced class of ML models currently corresponds to LLMs \citep{raiaan2024review} and LLM-based autonomous agents \citep{luo2025large}. 
Driven by scalable transformer architectures and aligned through reinforcement learning from human feedback, LLMs have moved beyond static text-to-text mapping to exhibit in-context learning \citep{dong2024survey}, semantic understanding, and zero-shot generalization capabilities \citep{kojima2022large}. 
The integration of these models into agentic frameworks has enabled a transition from passive generative systems to proactive decision-making entities. 
Such systems typically rely on modular architectures in which a central LLM is combined with explicit planning modules, memory structures (short- and long-term), and external tool-calling execution loops.

Despite these advances, the integration of ML into systems operating in social environments has raised increasing ethical concerns \citep{weidinger2021ethical, MLEthicsLaw, PhenotypesML}. 
These concerns stem from a dominant design paradigm that prioritizes technical performance metrics—such as predictive accuracy and computational efficiency \citep{regressionMetrics, rlMetrics, rlMetrics2, semi-supervisedMetrics, multi-class-metrics, ComputationalMetrics}—often at the expense of human-centered criteria.

Failure to incorporate such criteria into system design can lead to significant ethical issues \citep{electronics10050593, Human_metrics, ECThicGuidelines}. 
Key dimensions include sustainability, which concerns reducing energy consumption and improving model reusability to limit environmental impact \citep{strubell_etal_2019_energy}; transparency and explainability, which ensure that system decisions are interpretable and can be assessed by users \citep{ExplainabilityRequirements, causabilityMedicine, guidotti2018survey, XMLMetrics}; fairness, which addresses algorithmic bias and aims to prevent discriminatory outcomes across populations \citep{fairness, Onetofairness, AccuracyvsFairness}; and privacy and security, which focus on protecting data integrity against unauthorized access and adversarial attacks \citep{liu2020machine, Biggio_2018, SecML}. 
Finally, cultural competence refers to the ability of a system to adapt its behavior to the cultural norms and expectations of the environments in which it operates \citep{survey_fairness_wreprbias}.

Preliminary efforts have explored the integration of cultural competence into ML frameworks \citep{dlHumorRecognition, miehle2021multimodal, almaliki2016persuasive, khan2021one, dmytryk2022mining, karpouzis2024plato, Wu2012TheAO, dwivedi2024navigating, cirucci2024culturally, Quan2018, Zhou2022445, yun2024cic, mohamed2022artelingo, alharbi2023cultural, alharbi2024deep, jordan2023impact, castilla2022applying, yohe2023towards, blanchard2024cultural, bhatia2023gd, albuquerque2024hype, kuo2024identification}. 
However, this research direction remains in its early stages and is still limited, particularly in its application within the SR domain.

The first major limitation is that existing scholarship in this area focuses predominantly on data diversity \citep{not_uniform_data, seaborn2023not, lim2021social, hagerty2019global, stacy2025data, ijerph21121642, li2024health, lee2024negotiating, zhang2024ameliorating, blanchard2015socio, thakkar2024artificial}, empirically showing that contemporary datasets either omit or under-represent a wide range of cultural groups.

A second critical gap is the absence of standardized benchmark datasets specifically designed to capture cultural diversity. Such benchmarks are essential for rigorously evaluating the cultural competence of ML models within the SR domain \citep{Facial_recogn_eval_datasets, BiasDataset, not_uniform_data, seaborn2023not, lim2021social, hagerty2019global, stacy2025data, ijerph21121642, li2024health, lee2024negotiating, zhang2024ameliorating, blanchard2015socio, thakkar2024artificial}, in analogy with established evaluation frameworks used in algorithmic fairness research \citep{le2022survey}.

A third limitation concerns the absence of robust mitigation strategies for cultural incompetence in ML for SR. There is a lack of methodologies that can be consistently generalized and applied across large-scale ML environments.

It is important to note that the aforementioned studies—addressing social navigation \citep{patompak2017learning}, user clustering \citep{santos2018extended}, and gesture generation \citep{Gestures}—primarily operate under the implicit assumption of uniform data distributions across cultural groups. However, this assumption is often contradicted by empirical evidence. Subsequent research has highlighted significant disparities in data distribution, showing that minority cultural contexts are frequently under-represented or systematically excluded from foundational datasets \citep{not_uniform_data, seaborn2023not, lim2021social, hagerty2019global, stacy2025data, ijerph21121642, li2024health, lee2024negotiating, zhang2024ameliorating, blanchard2015socio, thakkar2024artificial}.

%\color{red}
\section{Cultural competence and robot functional domains}
%\color{black}
Table~\ref{ch1:tab:related} categorizes the functional domains in which cultural competence is relevant: Design, Verbal Interaction, Non-verbal Interaction, Social Distance and Navigation, Planning, Knowledge Representation, Generative Tasks, Activity Recognition, and Object Classification. The columns distinguish the methodological focus of the cited works: SR (implemented in SRs without ML), ML (generic ML techniques not specifically contextualized for SR applications), and SR+ML (ML frameworks specifically designed for SR applications).

As indicated in the table, this thesis represents a primary contribution to object classification and SR design using ML architectures specifically developed for SR applications, and it experimentally investigates the problem of embedding cultural competence using LLM-based architectures. The only other work identified within the first category pertains to a medical application, which highlights performance disparities that arise when models are exposed to culturally divergent datasets. These findings further support the need for robust mitigation strategies addressing cultural incompetence in ML and SR systems.

Within the aforementioned literature, some studies conceptualize cultural competence as a derivative of algorithmic fairness constraints \citep{karpouzis2024plato, albuquerque2024hype}. While a connection between these two fields exists, their objectives differ fundamentally: fairness constraints aim to decorrelate model outputs from sensitive attributes, whereas cultural competence aims to leverage cultural information as a feature to improve performance across cultural groups. Notwithstanding these differences, such studies are included to provide a comprehensive overview of the state of the art.

To the best of our knowledge, no prior research has explicitly addressed culturally competent object classification in the SR domain, nor has it experimentally investigated culture-as-a-nation approximations in this context.

%\begin{landscape}
\begin{table*}[t]
\label{ch1:tab:related}
\centering
\footnotesize
\setlength{\tabcolsep}{0.1cm}
%\renewcommand{\arraystretch}{1.0}
\begin{tabular}{l|c|c|c}
\toprule
\textbf{Category} & \textbf{SR} & \textbf{ML} & \textbf{SR+ML} \\
\midrule
Design & \makecell{\citep{akinade2023culturally, Dingjun2010}\\
\citep{ornelas2023redefining, Lee2016}\\
\citep{Papadopoulos_CARESSES, battistuzzi2021socially}\\
\citep{papadopoulos2020caresses, Shibata2009}\\
\citep{zantou2023culturally, 8673086}\\
\citep{vernon2024african, papadopoulos2022transcultural}\\
\citep{recchiuto2019designing}} 
& \makecell{\citep{miehle2021multimodal, khan2021one}\\
\citep{dmytryk2022mining, karpouzis2024plato}} 
& This study \\
\hline
Verbal Interaction & \makecell{\citep{Andrist2015, chugh2024aisles}\\
\citep{Bruno2018, nias2020culturebot}\\
\citep{bruno2016fuzzy, Rau2009}\\
\citep{obremski2022impact, CloudCultureDialogue}\\
\citep{recchiuto2020feasibility}} 
& \makecell{\citep{blanchard2024cultural}\\
\citep{bhatia2023gd}\\
\citep{almaliki2016persuasive, dlHumorRecognition}\\
\citep{albuquerque2024hype, Wu2012TheAO}\\
\citep{cirucci2024culturally, kuo2024identification}\\
\citep{dwivedi2024navigating}} 
& \citep{spitale2024appropriateness} \\
\hline
Non-verbal interaction & \makecell{\citep{Rehm2018, GreetingSelection}\\
\citep{pena2019eeva}} 
&  & \makecell{\citep{Gestures}\\
\citep{ren2025exploring}} \\
\hline
Social distance and navigation & \makecell{\citep{Eresha2013, Joosse2014}\\
\citep{Trovato2015, truong2016towards}} 
&  & \citep{patompak2017learning} \\
\hline
Planning & \citep{Planning, fedriga2023linguistic} &  &  \\
\hline
Knowledge representation & \makecell{\citep{GreetingSelection, KnowledgeRepresentation}\\
\citep{Grassi_2022, mansouri2020can}\\
\citep{bringsjord2020culturally}} &  &  \\
\hline
Generation &  & \makecell{\citep{Quan2018, Zhou2022445}\\
\citep{yun2024cic, mohamed2022artelingo}} &  \\
\hline
Object classification &  &  & This study \\
\hline
Other classification tasks & \citep{Chiang2019} 
& \makecell{\citep{alharbi2023cultural, alharbi2024deep}\\
\citep{jordan2023impact, castilla2022applying}\\
\citep{yohe2023towards}} 
& \citep{santos2018extended} \\
\bottomrule
\end{tabular}
\caption{Comparison of previous works exploring cultural competence. 
The rows represent specific aspects where cultural competence is relevant, while the columns indicate whether a related work has been implemented in social robots without relying on ML (SR), describes a generic ML technique not explicitly designed for social robots (ML), or introduces an ML technique explicitly designed for social robots (SR+ML).}
\label{ch1:tab:related}
\end{table*}
%\end{landscape}